% GENERATED FILE. Edit canonical lessons, then run npm run generate:books.
% canonical-source-hash: fnv1a64:43957b9ee1f31a85
% canonical-lessons: PT-C411-aluguer, PT-C411-condutor, PT-C411-matricula, PT-C411-quilometragem, PT-C411-franquia

\chapter{Um carro de aluguer}
\label{ch:pt-um-carro-de-aluguer}
\hlchaptermodality{411}

\begin{chapteropening}
\textbf{By the end of this chapter:} \emph{I can hire a car: the hire, the driver, the number plate, the mileage and the insurance excess.}

Complete the last lesson of chapter 411: I can hire a car: the hire, the driver, the number plate, the mileage and the insurance excess.
\end{chapteropening}

\section[o aluguer]{o aluguer — the hire, the rental}

\label{lesson:PT-C411-aluguer}



\begin{quote}
Before the new one: say the Portuguese for tarmac, then the Portuguese for to wind.
\end{quote}

\subsection*{You'll want to know: o aluguer}
\textbf{o aluguer} — \textquotedblleft{}the hire, the rental\textquotedblright{}.

\begin{quote}
O aluguer do carro custa quarenta euros por dia.
\end{quote}

\begin{quote}
\emph{The car hire costs forty euros a day.}
\end{quote}

In Brazil: \emph{o aluguel}. The verb is \emph{alugar}.

\begin{cousinweb}
\emph{aluguer} is inherited through Old Portuguese \emph{aloguer} from Latin \emph{allocāre} 'to place, to hire out' (EtymDB), built on \emph{locus} 'place'.
\end{cousinweb}

\begin{grammarlens}[title={What you've built}]
One more word for telling, explaining and describing at B1.
\end{grammarlens}

\subsection*{Your turn}
\begin{itemize}
\raggedright
  \item \emph{Say it:} \emph{o aluguer}
  \item \emph{Say it:} \emph{o aluguer}, once more
  \item \emph{Say it:} a sentence of your own with \emph{o aluguer}
  \item \emph{Recall:} say the Portuguese for tarmac, then the Portuguese for to wind, then say \emph{o aluguer} again
\end{itemize}

\begin{tcolorbox}[breakable,colback=teal!4,colframe=teal!35!black,title={Before you move on}]
What does o aluguer mean? (\textquotedblleft{}The hire, the rental\textquotedblright{}.) Say it once more.
\end{tcolorbox}
\section[o condutor]{o condutor — the driver}

\label{lesson:PT-C411-condutor}



\begin{quote}
Before the new one: say the Portuguese for to wind, then the Portuguese for the hire.
\end{quote}

\subsection*{You'll want to know: o condutor}
\textbf{o condutor} — \textquotedblleft{}the driver\textquotedblright{}.

\begin{quote}
O segundo condutor também tem de mostrar a carta.
\end{quote}

\begin{quote}
\emph{The second driver also has to show his licence.}
\end{quote}

\emph{A carta de condução} — the driving licence. \emph{O motorista} drives for a living.

\begin{cousinweb}
\emph{condutor} is a learned borrowing from Latin \emph{conductor} (EtymDB), from \emph{condūcere} 'to lead together', built on \emph{dūcere} 'to lead'.
\end{cousinweb}

\begin{grammarlens}[title={What you've built}]
One more word for telling, explaining and describing at B1.
\end{grammarlens}

\subsection*{Your turn}
\begin{itemize}
\raggedright
  \item \emph{Say it:} \emph{o condutor}
  \item \emph{Say it:} \emph{o condutor}, once more
  \item \emph{Say it:} a sentence of your own with \emph{o condutor}
  \item \emph{Recall:} say the Portuguese for to wind, then the Portuguese for the hire, then say \emph{o condutor} again
\end{itemize}

\begin{tcolorbox}[breakable,colback=teal!4,colframe=teal!35!black,title={Before you move on}]
What does o condutor mean? (\textquotedblleft{}The driver\textquotedblright{}.) Say it once more.
\end{tcolorbox}
\section[a matrícula]{a matrícula — the number plate; the registration; the enrolment}

\label{lesson:PT-C411-matricula}



\begin{quote}
Before the new one: say the Portuguese for the hire, then the Portuguese for the driver.
\end{quote}

\subsection*{You'll want to know: a matrícula}
\textbf{a matrícula} — \textquotedblleft{}the number plate; the registration; the enrolment\textquotedblright{}.

\begin{quote}
Um carro com matrícula estrangeira parou à porta.
\end{quote}

\begin{quote}
\emph{A car with foreign number plates stopped at the door.}
\end{quote}

At school or university, \emph{a matrícula} is enrolment.

\begin{cousinweb}
\emph{matrícula} is a learned borrowing from Late Latin \emph{mātrīcula} 'list, register' (EtymDB), a small form of \emph{mātrīx}, built on \emph{māter} 'mother'.
\end{cousinweb}

\begin{grammarlens}[title={What you've built}]
One more word for telling, explaining and describing at B1.
\end{grammarlens}

\subsection*{Your turn}
\begin{itemize}
\raggedright
  \item \emph{Say it:} \emph{a matrícula}
  \item \emph{Say it:} \emph{a matrícula}, once more
  \item \emph{Say it:} a sentence of your own with \emph{a matrícula}
  \item \emph{Recall:} say the Portuguese for the hire, then the Portuguese for the driver, then say \emph{a matrícula} again
\end{itemize}

\begin{tcolorbox}[breakable,colback=teal!4,colframe=teal!35!black,title={Before you move on}]
What does a matrícula mean? (\textquotedblleft{}The number plate; the registration; the enrolment\textquotedblright{}.) Say it once more.
\end{tcolorbox}
\section[a quilometragem]{a quilometragem — the mileage (in kilometres)}

\label{lesson:PT-C411-quilometragem}



\begin{quote}
Before the new one: say the Portuguese for the driver, then the Portuguese for the number plate.
\end{quote}

\subsection*{You'll want to know: a quilometragem}
\textbf{a quilometragem} — \textquotedblleft{}the mileage (in kilometres)\textquotedblright{}.

\begin{quote}
Neste aluguer, a quilometragem é ilimitada.
\end{quote}

\begin{quote}
\emph{With this hire, the mileage is unlimited.}
\end{quote}

Built on \emph{o quilómetro}.

\begin{cousinweb}
\emph{quilometragem} is built on \emph{quilómetro} with \emph{-agem}. \emph{Quilómetro} is a learned coinage from Greek \emph{khílioi} 'a thousand' and \emph{métron} 'measure'; EtymDB has no entry for \emph{quilómetro}.
\end{cousinweb}

\begin{grammarlens}[title={What you've built}]
One more word for telling, explaining and describing at B1.
\end{grammarlens}

\subsection*{Your turn}
\begin{itemize}
\raggedright
  \item \emph{Say it:} \emph{a quilometragem}
  \item \emph{Say it:} \emph{a quilometragem}, once more
  \item \emph{Say it:} a sentence of your own with \emph{a quilometragem}
  \item \emph{Recall:} say the Portuguese for the driver, then the Portuguese for the number plate, then say \emph{a quilometragem} again
\end{itemize}

\begin{tcolorbox}[breakable,colback=teal!4,colframe=teal!35!black,title={Before you move on}]
What does a quilometragem mean? (\textquotedblleft{}The mileage (in kilometres)\textquotedblright{}.) Say it once more.
\end{tcolorbox}
\section[a franquia]{a franquia — the excess (on insurance); the franchise}

\label{lesson:PT-C411-franquia}



\begin{quote}
Before the new one: say the Portuguese for the number plate, then the Portuguese for the mileage.
\end{quote}

\subsection*{You'll want to know: a franquia}
\textbf{a franquia} — \textquotedblleft{}the excess (on insurance); the franchise\textquotedblright{}.

\begin{quote}
Com este seguro, a franquia é de trezentos euros.
\end{quote}

\begin{quote}
\emph{With this insurance, the excess is three hundred euros.}
\end{quote}

The excess is the part you pay yourself after an accident.

\begin{cousinweb}
\emph{franquia} is built on \emph{franco} 'free, exempt', which was borrowed from French \emph{franc} (EtymDB), first the name of the Franks: a \emph{franquia} was a freedom from a charge.
\end{cousinweb}

\begin{grammarlens}[title={What you've built}]
One more word for telling, explaining and describing at B1.
\end{grammarlens}

\subsection*{Your turn}
\begin{itemize}
\raggedright
  \item \emph{Say it:} \emph{a franquia}
  \item \emph{Say it:} \emph{a franquia}, once more
  \item \emph{Say it:} a sentence of your own with \emph{a franquia}
  \item \emph{Recall:} say the Portuguese for the number plate, then the Portuguese for the mileage, then say \emph{a franquia} again
\end{itemize}

\begin{tcolorbox}[breakable,colback=teal!4,colframe=teal!35!black,title={Before you move on}]
What does a franquia mean? (\textquotedblleft{}The excess (on insurance); the franchise\textquotedblright{}.) Say it once more.
\end{tcolorbox}
